% Report fragment. Requires amsmath, amssymb, amsthm and a theorem/lemma environment.
\section{PA194N: a ten-entry no-push proof-search controller}
\label{sec:np-controller}

The primary PA194 table is modified at exactly one native transition:
\[
 (164,1):\quad (1,R,98)\quad\longmapsto\quad(1,R,114).
\]
The seven axiom-reader folds and the complete 190-token separated-head
emitter remain unchanged. The result has 194 working states and 388 rows.
Its source has 3,654 primitive instruction sites instead of 3,673.
This is a change of command semantics: successful modus ponens no longer
pushes a zero onto the proof stack. It retains the original Shift test and
retains the conclusion in the current top, distinguishing it from the
cached-head PA194C and symmetric-guard PA194CS controllers.

The no-push command rule itself was proposed earlier in
\texttt{conpa\_tm\_comparison\_and\_reduction\_search}, at 196 states.
Its ZFC ancestry is also explicit in
\texttt{candidate-k215/alternatives/zero-spare-transfer/report.tex},
whose ``Removing the successful MP zero push section gives the same
store update and reserved-zero proof compiler for a different Shift guard.
The zero-alias branch-coalescence principle is inherited, not claimed
as a new general idea.
The new contribution here is its complete 194-state realization with the
separated-head emitter, the full symbolic primitive-contract certificate,
and the specialized proof compiler below. The command proof is an ordinary
mathematical proof. The Python identities and tape certificates are not an
end-to-end Lean or Rocq proof of the consistency equivalence.

\subsection{Exact numerical operations and command contracts}

Let
\[
 p(a,b)=\frac{(a+b)(a+b+1)}2+b,\qquad [a;w]=p(a,w),\qquad
 I(a,b)=4p(a,b),\qquad U(i,a)=4p(a,i)+2.
\]
Iterated brackets put their leftmost entry at the top. The Cantor
projections are total and satisfy
\[
 \operatorname{hd}[a;w]=a,\quad \operatorname{tl}[a;w]=w,
 \quad[0;0]=0,\quad\operatorname{hd}(0)=\operatorname{tl}(0)=0.
\]
A numerical stack is therefore not being treated as a free list.
A command store is $(t,w,d,h)$: current formula, stack, offset, and cached
formula. Code zero is truth and code one is the target contradiction.
The new MP command, with $y=d+x$, is
\begin{equation}\label{eq:np-mp}
 \mathsf{MP}_x(t,w,d,h)=
 \begin{cases}
 (y,w,y,0),& I(t,y)=h,\\
 (0,w,y,0),& I(t,y)\ne h.
 \end{cases}
\end{equation}
Shift is retained literally:
\begin{equation}\label{eq:np-shift}
 S(t,w,d,h)=
 \begin{cases}
 \mathsf{halt},&t+h=1,\\
 (h,\operatorname{tl}(w),d,\operatorname{hd}(w)),&t+h\ne1.
 \end{cases}
\end{equation}
For $D=d+x$, $H=h+y$, the unchanged AX command has output
\begin{equation}\label{eq:np-ax}
 \mathsf{AX}_{A,x,y}(t,w,d,h)
   =\bigl(g(H,D,A),[q_1,\ldots,q_9,t;w],0,0\bigr).
\end{equation}
Here $q_1=\Gamma(A)=I(A,U(A,A))$, $q_2,\ldots,q_7$ are the
six retained PA196 heads in reverse emission order,
\[
 q_8=\rho(D,A,H,d)=I(I(I(H,d),D),I(I(D,H),I(A,H))),
 \qquad q_9=U(d,h).
\]
The gate $g$ is the original separate-head guarded separation gate.
The free-constructor computation in
\texttt{check\_emitter\_and\_words.py} binds all these formula trees to
the delivered grammar. All $q_1,\ldots,q_8,g$ are divisible by four,
and $q_9\equiv2\pmod4$, with $q_9\ge2$. In particular none is code one.
This is a statement about constructor values, requiring no consistency
assumption and holding for arbitrary incoming natural stores.

The count in (\ref{eq:np-ax}) is nine generated stack entries plus the
saved old top. The two identities needed below are
\begin{align}
 \mathsf{AX}_{A,x,y};S^9(t,w,d,h)
   &=(q_8,[t;w],0,U(d,h)),\label{eq:np-nine}\\
 \mathsf{AX}_{A,x,y};S^{10}(t,w,d,h)
   &=(U(d,h),w,0,t).\label{eq:np-ten}
\end{align}
All displayed Shifts are nonhalting: their tests involve two generated
non-one codes, zero and a generated code, or the universal code of value
at least two. This also establishes exact stack effects when zero pairs
collapse.

\subsection{Uniform stack operations and a one-command discard}

Write $A_0=\mathsf{AX}_{0,0,0}$ and define
\[
 \mathsf{Push}=A_0;S^9;\mathsf{MP}_0,\qquad
 \mathsf{Pop}=A_0;S^{11},\qquad
 \mathsf{Discard}=\mathsf{MP}_0,\qquad
 \mathsf{Offset}_i=\mathsf{MP}_i.
\]
Their exact contracts are
\begin{align}
 \mathsf{Push}(t,w,d,h)&=(0,[t;w],0,0),\label{eq:np-push}\\
 \mathsf{Pop}(t,[a;w],d,h)&=(t,w,0,a),\label{eq:np-pop}\\
 \mathsf{Discard}(t,w,0,h)&=(0,w,0,0),\label{eq:np-discard}\\
 \mathsf{Offset}_i(0,w,0,0)&=(0,w,i,0).\label{eq:np-offset}
\end{align}
For Push, (\ref{eq:np-nine}) makes the final MP compare an implication
with a universal code, so it fails with offset zero. For Pop, the first
ten Shifts give (\ref{eq:np-ten}). The eleventh Shift tests
$U(d,h)+t\ge2$, hence it does not halt even when $t=1$.
It pops exactly the displayed item. Thus Pop is total on all natural
stores. In particular $\mathsf{Pop}(t,0,d,h)=(t,0,0,0)$ and
$\mathsf{Push}(0,0,d,h)=(0,0,0,0)$.

For Discard, the two MP outcomes in (\ref{eq:np-mp}) coincide when
$d=x=0$, for arbitrary $h$: both are $(0,w,0,0)$. This gives a single-command discard,
where the old dummy-push controller needed a longer word. The old cache
need not be zero for this identity. For Offset,
if $i>0$ then $I(0,i)>0$, so the comparison fails; if $i=0$ it succeeds
and returns $(0,w,0,0)$. The same one-command definition covers both cases.
These simplifications affect the proof-word compiler; they do not claim
additional physical state savings.

\subsection{Axiom selectors and generalization}

In the following contracts the arbitrary incoming store is $(t,w,0,0)$.
For $1\le j\le7$, the word
\[
 \mathsf{Discard};\mathsf{AX}_{A,D,H};S^{j+1};\mathsf{Pop}^{9-j}
\]
returns $(q_j,w,0,0)$ with the parameters evaluated at the desired
$A,D,H$. After the initial Shifts, the remaining $9-j$ stack entries
are exactly the unselected generated tail and the saved zero. The Pops
preserve the selected top and consume those entries; the last pop puts
that saved zero into the cache. For the dominated head, if $u\le a$,
use
\[
 \mathsf{Discard};\mathsf{Offset}_u;
 \mathsf{AX}_{b,a-u,c};S^9;\mathsf{Pop}.
\]
It returns $(\rho(a,b,c,u),w,0,0)$, with the numerical dominance condition
coming from $a=u+(a-u)$. A permitted gate is obtained by
\[
 \mathsf{Discard};\mathsf{AX}_{A,D,H};\mathsf{Pop}^{10},
\]
which returns $(g(H,D,A),w,0,0)$. Truth uses Discard alone.

The word
\[
 \mathsf{Gen}_i=\mathsf{Push};\mathsf{Offset}_i;S;
                \mathsf{AX}_{i,0,0};S^{10}
\]
returns $(U(i,a),w,0,0)$ from $(a,w,0,0)$. Before its AX the store is
$(0,w,i,a)$; its last ten Shifts use (\ref{eq:np-ten}) with saved top zero.
Its first Shift tests $0+0$, and the later ones are among the safe AX
shifts. Hence this generalization word is total even for $a=1$.

\subsection{Modus ponens with its spare zero supplied by the compiler}

Suppose recursively compiled proofs return $a$ and $I(a,b)$ with the
stack framed and both caches cleared, or halt earlier. The MP compiler is
\begin{equation}\label{eq:np-compile}
 \mathsf{Discard};\mathsf{Push};C(I(a,b));\mathsf{Push};
 C(a);\mathsf{Push};S;S;\mathsf{MP}_b;\mathsf{Pop}.
\end{equation}
The first two commands reserve a zero below the premise computations.
If the recursive calls have returned, the store before the two Shifts is
\[
             (0,[a,I(a,b),0;w],0,0).
\]
The first Shift returns $(0,[I(a,b),0;w],0,a)$.
The second halts if $a=1$; otherwise it returns
$(a,[0;w],0,I(a,b))$. MP succeeds and returns
\[
                         (b,[0;w],b,0).
\]
By the total Pop contract, the final word returns $(b,w,0,0)$, including
$b=1$. The stack zero is now consumed once, at the compiler level.
At $a=b=w=0$ the same proof uses the literal Cantor projection identities;
no extra distinction between a zero push and an empty stack is introduced.

Induction on a finite derivation $p$ in the inherited separate-head
calculus therefore proves:
\begin{quote}
From every $(t,w,0,0)$, $C(p)$ either halts or returns
$(\operatorname{concl}(p),w,0,0)$.
\end{quote}
The axiom and generalization cases use the words above. The MP case is
(\ref{eq:np-compile}); extending a word after an early halt still halts.
In particular a refutation compiler ends in top one, if it has not already
halted. Its final appended Shift halts because $1+0=1$.
For arbitrary incoming $t,w,d,h$, first apply $A_0$ to clear $d,h$.
Thus the concrete word
\[
                         A_0;C(p);S
\]
halts from \emph{every} incoming store whenever $p$ derives code one.
The compiler is a primitive-recursive traversal of the finite derivation,
and PA proves its correctness by induction on that coded derivation.
It does not assume a reset between the scheduler's unsuccessful trials.

The unchanged scheduler enumerates finite command words with the usual
Cantor continuation coding. A cofinal family of halting fresh candidates
is obtained by prefixing the universal word by
$\mathsf{AX}_{k,0,0}$. Let $E(\epsilon)=0$ and let a positive continuation be encoded by
$E(c::w)=1+\operatorname{encodeCommand}(c,E(w))$. In particular
$\operatorname{encodeCommand}(\mathsf{AX}_{A,x,y},z)
=p(A+2,p(y,p(x,z)))$. Put $z=E(A_0;C(p);S)$, using this positive
continuation convention rather than a fresh-candidate code. The fresh
candidate codes for the prefixed words are exactly
$p(k+2,p(0,p(0,z)))$, which are strictly increasing and
unbounded in $k$. The fresh counter begins at two, so omitting smaller
candidates loses none of this cofinal family. Every decoded continuation
strictly decreases, and the unchanged source macro loops terminate.
The resulting finite-word and finite-trace construction is formalizable
in PA by primitive recursion and induction.

\subsection{Soundness and the PA calibration}

Annotate the current top $t$, cached head $h$, and every live canonical
Cantor-stack entry by a derivation in the separate-head calculus, or by a
PA proof of its arithmetic translation. For $w>0$, its decoded tail is
strictly smaller than $w$; thus the canonical stack and its annotation
have a primitive-recursive bounded construction. A zero stack supplies
truth when popped, and pushing zero onto zero preserves the empty
canonical annotation.

AX adds the unchanged valid axiom instances, a generalization of the old
cached theorem, and the saved old top. On a successful MP, the annotations
of $t$ and $I(t,y)=h$ produce an annotation for $y$. The stack is retained,
and the cache receives the truth proof. Failure outputs truth and retains
the stack. Shift transports the annotations. If its test succeeds, then
$t+h=1$ implies that one of the two annotated formulas is code one.
The inherited syntactic arithmetic translation turns that annotation into
a PA refutation. No annotation is imposed on the arithmetic offset $d$.

Consequently PA proves both finite syntactic reductions:
\[
 \text{source halt}\ \longrightarrow\ \text{PA refutation},
 \qquad
 \text{PA refutation}\ \longrightarrow\ \text{source halt}.
\]
For the second reduction, compose the supplied primitive-recursive
PA-to-separated-head proof translation with the all-store compiler above.
For the first, compose the proof annotation with the supplied
separated-head-to-PA translation. This uses the existing logical calibration,
including its treatment of the code-zero alias and coincident variable
labels. Nothing in the no-push modification changes that calculus or
requires a new logical axiom, a reflection principle, or consistency as
an additional assumption.

\subsection{Exact source splice and its physical realization}

At a successful MP comparison, the source has $T=H=X=Y=0$, $D=b$, and
arbitrary stack $W=w$. The original subtree is
\[
 \texttt{node163}=(\texttt{node6},\texttt{node55}),\qquad
 \texttt{node55}=(\texttt{node49},\texttt{node54}).
\]
The first child of node55 pushes the current zero; node54 copies $D$ to
$T$ while preserving $D$. PA194N replaces node163 by
$(\texttt{node6},\texttt{node54})$.
The deleted source interval is $[3634,3653)$ in the original zero-based
numbering. It contains 19 instructions and has exactly one exterior
entry, the zero branch of instruction 3628. Every remaining instruction,
including every successor, equals the old instruction under the explicit
index relocation in \texttt{build\_nopush.py}, with that entry redirected
to the retained copy loop. Both failed MP and every other command retain
their source effects.

The retained copy loop, in local labels, is
\[
\begin{array}{c|l}
0&\mathrm{DEC}\ D\quad(0\mapsto4,\ >0\mapsto1)\\
1&\mathrm{INC}\ T;\ \mathrm{goto}\ 2\\
2&\mathrm{INC}\ X;\ \mathrm{goto}\ 3\\
3&\mathrm{goto}\ 0\\
4&\mathrm{DEC}\ X\quad(0\mapsto7,\ >0\mapsto5)\\
5&\mathrm{INC}\ D;\ \mathrm{goto}\ 6\\
6&\mathrm{goto}\ 4
\end{array}
\]
Label seven is the exit. After $j\le b$ iterations of the first loop,
$(T,D,X)=(j,b-j,j)$. After its final zero test and $j\le b$ iterations
of the second, $(T,D,X)=(b,j,b-j)$. Hence the exit is
$(T,D,X)=(b,b,0)$, with all other registers unchanged, in exactly
$7b+2$ primitive source steps. These two elementary induction invariants
prove the retained loop's contract over all naturals.

The literal physical table is rebuilt from the pinned grammar and seven
folds. The changed native transition enters the already existing reader
for node54 directly. All 4,755 arbitrary-register primitive outcomes pass
(the maximum symbolic-checker step count is 497); all 28 local fold
contexts pass simultaneously; the 23 arithmetic-engine controls retain
every original row. The complete blank-tape initialization reaches native
control 84 at head $-2$, with the entire clean frame and 26 spare blocks,
after exactly 58,154 literal transitions. The startup equality includes
every nonzero tape cell.

The generic register-block and unary-scan arguments are elementary
inductions over coded finite lists and natural lengths. Their trace
witnesses are primitive recursively constructible. The finite successful
primitive certificates therefore give PA-provable positive finite source
simulation. Conversely, given a physical first halt at time $T$, iterate the totalized
source step function for $T+1$ instructions. If this bounded iteration
finds a source halt, reflection is established. Otherwise its positive
nonhalting simulations, preceded by startup, produce a deterministic
physical trace beyond that first halt, a contradiction. This proves
reflection by a bounded construction.
Combining source calibration and physical simulation yields the ordinary
mathematical conclusion
\[
 \mathrm{PA}\vdash
 \bigl(\neg\operatorname{Halt}(\mathrm{PA194N})
                 \leftrightarrow\operatorname{Con}(\mathrm{PA})\bigr).
\]
The finite and symbolic executables support this argument and bind it to
the exact table. They do not replace the explicitly identified inherited
logical proof transformations, and no newly compiled end-to-end trusted
kernel endpoint is claimed.
